% GENERATED FILE. Edit canonical lessons, then run npm run generate:books.
% canonical-source-hash: fnv1a64:a54f8ed82044d3de
% canonical-lessons: KA-C180-ucita, KA-C180-sarkari, KA-C180-khasagi, KA-C180-sthaliya, KA-C180-kaddaya

\chapter{Free of charge, Government, Private, Local, Compulsory}
\label{ch:ka-free-of-charge-government-private-local-compulsory}
\hlchaptermodality{180}

\begin{chapteropening}
\textbf{By the end of this chapter:} \emph{I can say free of charge, government, private, local and compulsory.}

Complete the last lesson of chapter 180: I can say free of charge, government, private, local and compulsory.
\end{chapteropening}

\section[\texorpdfstring{\kn{ಉಚಿತ} (ucita)}{ಉಚಿತ (ucita)}]{\kn{ಉಚಿತ} (ucita) — free of charge}

\label{lesson:KA-C180-ucita}



\begin{quote}
Before the new one: say the Kannada for bland, then the Kannada for expensive.
\end{quote}

\subsection*{You'll want to know: \kn{ಉಚಿತ}}
\textbf{\kn{ಉಚಿತ}} — \emph{ucita} — \textquotedblleft{}free of charge\textquotedblright{}.

\begin{grammarlens}[title={What you've built}]
One more word for this chapter.
\end{grammarlens}

\subsection*{Your turn}
\begin{itemize}
\raggedright
  \item \emph{Say it:} \emph{ucita}
  \item \emph{Say it:} \emph{ucita}, once more
  \item \emph{Say it:} say \emph{ucita}
  \item \emph{Recall:} say the Kannada for bland, then the Kannada for expensive, then say \emph{ucita} again
\end{itemize}

\begin{tcolorbox}[breakable,colback=teal!4,colframe=teal!35!black,title={Before you move on}]
What does \kn{ಉಚಿತ} mean? (\textquotedblleft{}Free of charge\textquotedblright{}.) Say it once more.
\end{tcolorbox}
\section[\texorpdfstring{\kn{ಸರ್ಕಾರಿ} (sarkāri)}{ಸರ್ಕಾರಿ (sarkāri)}]{\kn{ಸರ್ಕಾರಿ} (sarkāri) — government, governmental}

\label{lesson:KA-C180-sarkari}



\begin{quote}
Before the new one: say the Kannada for expensive, then the Kannada for free of charge.
\end{quote}

\subsection*{You'll want to know: \kn{ಸರ್ಕಾರಿ}}
\textbf{\kn{ಸರ್ಕಾರಿ}} — \emph{sarkāri} — \textquotedblleft{}government, governmental\textquotedblright{}.

\begin{grammarlens}[title={What you've built}]
One more word for this chapter.
\end{grammarlens}

\subsection*{Your turn}
\begin{itemize}
\raggedright
  \item \emph{Say it:} \emph{sarkāri}
  \item \emph{Say it:} \emph{sarkāri}, once more
  \item \emph{Say it:} say \emph{sarkāri}
  \item \emph{Recall:} say the Kannada for expensive, then the Kannada for free of charge, then say \emph{sarkāri} again
\end{itemize}

\begin{tcolorbox}[breakable,colback=teal!4,colframe=teal!35!black,title={Before you move on}]
What does \kn{ಸರ್ಕಾರಿ} mean? (\textquotedblleft{}Government, governmental\textquotedblright{}.) Say it once more.
\end{tcolorbox}
\section[\texorpdfstring{\kn{ಖಾಸಗಿ} (khāsagi)}{ಖಾಸಗಿ (khāsagi)}]{\kn{ಖಾಸಗಿ} (khāsagi) — private}

\label{lesson:KA-C180-khasagi}



\begin{quote}
Before the new one: say the Kannada for free of charge, then the Kannada for government.
\end{quote}

\subsection*{You'll want to know: \kn{ಖಾಸಗಿ}}
\textbf{\kn{ಖಾಸಗಿ}} — \emph{khāsagi} — \textquotedblleft{}private\textquotedblright{}.

\begin{grammarlens}[title={What you've built}]
One more word for this chapter.
\end{grammarlens}

\subsection*{Your turn}
\begin{itemize}
\raggedright
  \item \emph{Say it:} \emph{khāsagi}
  \item \emph{Say it:} \emph{khāsagi}, once more
  \item \emph{Say it:} say \emph{khāsagi}
  \item \emph{Recall:} say the Kannada for free of charge, then the Kannada for government, then say \emph{khāsagi} again
\end{itemize}

\begin{tcolorbox}[breakable,colback=teal!4,colframe=teal!35!black,title={Before you move on}]
What does \kn{ಖಾಸಗಿ} mean? (\textquotedblleft{}Private\textquotedblright{}.) Say it once more.
\end{tcolorbox}
\section[\texorpdfstring{\kn{ಸ್ಥಳೀಯ} (sthaḷīya)}{ಸ್ಥಳೀಯ (sthaḷīya)}]{\kn{ಸ್ಥಳೀಯ} (sthaḷīya) — local}

\label{lesson:KA-C180-sthaliya}



\begin{quote}
Before the new one: say the Kannada for government, then the Kannada for private.
\end{quote}

\subsection*{You'll want to know: \kn{ಸ್ಥಳೀಯ}}
\textbf{\kn{ಸ್ಥಳೀಯ}} — \emph{sthaḷīya} — \textquotedblleft{}local\textquotedblright{}.

\begin{grammarlens}[title={What you've built}]
One more word for this chapter.
\end{grammarlens}

\subsection*{Your turn}
\begin{itemize}
\raggedright
  \item \emph{Say it:} \emph{sthaḷīya}
  \item \emph{Say it:} \emph{sthaḷīya}, once more
  \item \emph{Say it:} say \emph{sthaḷīya}
  \item \emph{Recall:} say the Kannada for government, then the Kannada for private, then say \emph{sthaḷīya} again
\end{itemize}

\begin{tcolorbox}[breakable,colback=teal!4,colframe=teal!35!black,title={Before you move on}]
What does \kn{ಸ್ಥಳೀಯ} mean? (\textquotedblleft{}Local\textquotedblright{}.) Say it once more.
\end{tcolorbox}
\section[\texorpdfstring{\kn{ಕಡ್ಡಾಯ} (kaḍḍāya)}{ಕಡ್ಡಾಯ (kaḍḍāya)}]{\kn{ಕಡ್ಡಾಯ} (kaḍḍāya) — compulsory}

\label{lesson:KA-C180-kaddaya}



\begin{quote}
Before the new one: say the Kannada for private, then the Kannada for local.
\end{quote}

\subsection*{You'll want to know: \kn{ಕಡ್ಡಾಯ}}
\textbf{\kn{ಕಡ್ಡಾಯ}} — \emph{kaḍḍāya} — \textquotedblleft{}compulsory\textquotedblright{}.

\begin{grammarlens}[title={What you've built}]
One more word for this chapter.
\end{grammarlens}

\subsection*{Your turn}
\begin{itemize}
\raggedright
  \item \emph{Say it:} \emph{kaḍḍāya}
  \item \emph{Say it:} \emph{kaḍḍāya}, once more
  \item \emph{Say it:} say \emph{kaḍḍāya}
  \item \emph{Recall:} say the Kannada for private, then the Kannada for local, then say \emph{kaḍḍāya} again
\end{itemize}

\begin{tcolorbox}[breakable,colback=teal!4,colframe=teal!35!black,title={Before you move on}]
What does \kn{ಕಡ್ಡಾಯ} mean? (\textquotedblleft{}Compulsory\textquotedblright{}.) Say it once more.
\end{tcolorbox}
